\section{结论}

\subsection{问题二}

在现有 21 个催化剂组合中，250--350 ℃的乙醇转化率和 $\mathrm{C}_4$ 烯烃选择性端点变化均为正，说明温度端点促进方向在已实验范围内最一致。严格匹配与跨背景比较表明，Co 负载量由 1 wt\% 提高至 5 wt\% 时，其对乙醇转化率的影响可随匹配背景反转，而该具体水平变化对选择性的平均差在两个可复现背景中均为负；乙醇进料条件和装料方式均不存在脱离背景与温度的统一优劣。固定总装料量时，组分质量分配对两个响应的方向不同；总装料量也不呈现越大越优的规律。

在当前五因素定义和严格匹配条件下，仅 $(m_{\mathrm{cat}},q_{\mathrm{EtOH}})$ 与 $(q_{\mathrm{EtOH}},z)$ 两类完整局部四格可以直接识别百分点加法尺度的二阶非加性。该结论仅适用于所比较的离散水平和实际共同温度，不延伸为全局交互强度、因素重要性排序或严格因果判断。
